\documentclass[11pt]{article}
\usepackage{mathblatt}
\usepackage{multicol}
% Probe Serie 08.10.2026: Übersichtsblatt Kurvenuntersuchung, Fassung 2 (Stichpunkte), von Hand gesetzt.
% Reihenfolge und Rand/Kern aus Fassung 1 (bau/proben/2026-10-08/uebersicht-kurvenuntersuchung/);
% Handgriffe aus mathe-nachhilfe abitur/gliederung/kurvenuntersuchung.md (Zuschnitt-Zeilen, 50 Sorten);
% "selten" = Sorte mit höchstens einer Aufgabe 2022–2026. Keine Aufgaben, keine Lösungen.
\newcommand{\kennung}{R4V}
\newcommand{\ueblock}[1]{\par\addvspace{10pt}\noindent{\small\scshape\color{mbgrau}#1}\par\nobreak\vspace{-2pt}%
  \noindent{\color{mbgitter}\rule{\linewidth}{0.4pt}}\par\nobreak\vspace{1pt}}
\newcommand{\ueab}[2]{\par\addvspace{6pt}\noindent{\bfseries #1}\ifblank{#2}{}{\hspace{0.6em}{\footnotesize\color{mbgrau}Rand}}\par\nobreak\vspace{2pt}}
\newcommand{\uf}[1]{\par\noindent\hangindent=1.25em\hangafter=1\makebox[1.25em][l]{$\square$}{\small #1}\par\vspace{3.5pt}}
\newcommand{\kl}[1]{{\footnotesize\color{black!70}(#1)}}
\newcommand{\selten}{\hspace{0.4em}{\footnotesize\color{mbgrau}selten}}
\setlength{\columnsep}{9mm}

\begin{document}
\pfheftstil
\fancyfoot[C]{{\footnotesize\color{mbgrau}\kennung\hfill\thepage}}
\pfheftkopf{Kurvenuntersuchung}{Abi GK \textperiodcentered{} Übersicht}
\noindent Was ist gemeint? Kannst du das? Frag, wenn du unsicher bist.
\vspace{2pt}

\begin{multicols}{2}
\raggedright\setlength{\parskip}{0pt}
\ueblock{Eigenschaften am Term}

\ueab{Graphen verschieben und spiegeln}{}
\uf{$-f(x)$, $f(-x)$: An welcher Achse gespiegelt? \kl{Drehung, selten}}
\uf{$f(x-2)$, $f(x)+3$, $2f(x)$: Was macht der Graph?}
\uf{Wie wird der eine Graph zum anderen?\selten}

\ueab{Für große und kleine $x$}{}
\uf{Wohin für $x\to\pm\infty$? Höchste Potenz ansehen.}
\uf{Und bei $(4-x)\cdot e^{x}$? Wer gewinnt?}

\ueab{Symmetrie}{rand}
\uf{Symmetrisch? Sieh auf die Exponenten. \kl{mit $f(-x)$, selten}}

\ueab{Nullstellen und Achsenschnittpunkte}{}
\uf{Achsenschnittpunkte: $f(0)$ und $f(x)=0$, auch mit $e^{x}$.}
\uf{Ausklammern; $e^{x}$ ist nie null. \kl{Substitution}\selten}
\uf{Wie viele Nullstellen? Am Produkt ablesen.\selten}

\ueblock{Ableitung}

\ueab{Wo steigt, wo fällt der Graph?}{}
\uf{Am Graphen von $f'$: über oder unter der $x$-Achse?}
\uf{Am Term von $f'$ oder aus Hoch- und Tiefpunkt.\selten}
\uf{Was sagen $f(a)$ und $f'(a)$ über den Graphen?\selten}

\ueab{Hoch- und Tiefpunkte}{}
\uf{Alle Extrempunkte: $f'(x)=0$, dann $f''$ prüfen.}
\uf{Extrempunkt bei $x_0$? Zeige $f'(x_0)=0$.}
\uf{Bei $(x+3)\cdot e^{-x}$: Produktregel. \kl{Verbindungsgerade}}

\ueab{Wendepunkte}{}
\uf{Wendepunkt: $f''(x)=0$ -- oder am Graphen ablesen.}
\uf{Wendepunkt bei $e$-Funktion: zweimal ableiten.}
\uf{Was bedeutet der Wendepunkt in der Sache?\selten}

\ueab{Den Graphen skizzieren}{rand}
\uf{Graph aus Hoch-, Tief-, Wendepunkt skizzieren.\selten}
\uf{Werte ablesen, Verlauf in der Sache beschreiben.\selten}

\ueblock{Anwenden}

\ueab{Maße aus dem Graphen}{}
\uf{Wie hoch, wie breit? Werte mal Maßstab. \kl{Kreisbögen}}
\uf{Welcher Quader passt um das Profil?}

\ueab{Größte Fläche, größter Abstand}{rand}
\uf{Rechteck unter dem Graphen: $A(u)$ maximieren.\selten}
\uf{Abstand $d(x)=f(x)-g(x)$ maximieren.\selten}

\ueab{Gleichungen und Ungleichungen}{}
\uf{Schnittstelle von $f$ und $g$ nachweisen.\selten}
\uf{Wo ist $f>g$? Differenz faktorisieren. \kl{Betrag}\selten}
\uf{Grafisch lösen; $f'(x)=$ mittlere Steigung.\selten}

\ueab{Funktionsgleichung aus Bedingungen}{rand}
\uf{Gerade durch zwei Punkte aufstellen.\selten}
\uf{Gleichungen für $a,b,c,d$ aufstellen. \kl{knickfrei}\selten}

\vfill
\noindent{\small\color{mbgrau}Weiter: Graphen verschieben und spiegeln}
\end{multicols}
\end{document}
